% ============================================================
%  medios_de_contraste.tex --- SOLO CONTENIDO
%  No tiene \documentclass ni \begin{document}.
%  Se incluye desde main.tex con \input{medios_de_contraste}
% ============================================================

\chapter{Medios de Contraste}

% --- Datos del docente ---
{\normalsize\color{gray} Lic. Matías Prado \quad|\quad Lic. Mauro Canclini}
\vspace{1em}

% ============================================================
%  PUNTOS CLAVE PARA EL EXAMEN ORAL
% ============================================================
\begin{cajakeypoints}
\textbf{\color{azulprincipal} Puntos clave para el examen oral}
\begin{itemize}[noitemsep, topsep=2pt]
  \item Definición y finalidad diagnóstica/terapéutica de los MC.
  \item Triple clasificación: por tipo de imagen, vía de administración y método de imagen.
  \item Estructura química de los MCI: anillo de benceno, posición de los átomos de yodo.
  \item Propiedades fisicoquímicas: osmolaridad, ionicidad, estructura molecular, viscosidad.
  \item El contraste \textbf{ideal}: no iónico, dímero, isoosmolar $\Rightarrow$ menor incidencia de reacciones adversas.
  \item Reacciones adversas: fisicoquímicas (dosis-dependientes) vs.\ idiosincráticas (dosis-independientes).
  \item Clasificación por gravedad: leves (98\%), moderadas (1\%), graves (1\%).
  \item NIC: filtrado glomerular (FG) como predictor; umbrales de riesgo.
  \item Protocolo de extravasación: pasos ordenados.
  \item Gadolinio: paramagnético, quelante, fibrosis sistémica nefrogénica en insuficiencia renal.
  \item Consentimiento informado: deber médico y derecho del paciente.
  \item Situaciones especiales: embarazo y lactancia.
\end{itemize}
\end{cajakeypoints}

\vspace{1em}

% ============================================================
\section*{Definición}
% ============================================================

\begin{cajadefinicion}
Un \textbf{medio de contraste (MC)} es toda sustancia o combinación de sustancias que, introducida en el organismo por cualquier vía, \textbf{permite resaltar y opacificar estructuras anatómicas normales} (órganos, vasos) \textbf{y patológicas} (tumores). También evalúa la perfusión y diferencia interfases de densidad entre tejidos de forma \textit{temporal}, con fines médicos diagnósticos y/o terapéuticos.
\end{cajadefinicion}

La indicación del MC se justifica porque las \textbf{partes blandas poseen bajo contraste natural}, lo que hace necesaria la administración de sustancias que las diferencien. El MC se elige y adapta según la región a estudiar y la técnica de exploración.

% ============================================================
\section*{Clasificación}
% ============================================================

\subsection*{1. Por tipo de imagen}

\renewcommand{\arraystretch}{1.35}
\begin{center}
\begin{tabular}{>{\bfseries}p{2.8cm} p{5cm} p{5cm}}
\hline
Tipo & Mecanismo & Ejemplos \\
\hline
Positivos & Atenúan más los Rx que los tejidos blandos. Se ven \textbf{radiopacos} (blancos). & Bario, Yodo, Gadolinio$^*$ \\[0.5em]
\hline
Negativos & Atenúan menos los Rx que los tejidos blandos. Se ven \textbf{radiolúcidos} (negros). & Aire, CO$_2$ \\[0.5em]
\hline
Neutros & Distienden y rellenan el tubo digestivo sin modificar densidad. & Agua, Metilcelulosa, Polietilenglicol, Manitol \\
\hline
\multicolumn{3}{l}{\footnotesize $^*$El gadolinio no interactúa con los Rx; se emplea en resonancia magnética.}
\end{tabular}
\end{center}

\subsection*{2. Por vía de administración}

Intravenosa (EV), intraarterial, intracanalicular, intraarticular, intravaginal, rectal, oral, entre otras.

\subsection*{3. Por método de imagen}

\renewcommand{\arraystretch}{1.35}
\begin{center}
\begin{tabular}{>{\bfseries}p{2.5cm} p{9.5cm}}
\hline
Método & Tipo de contraste utilizado \\
\hline
Rx contrastada & Bario, aire/polvo efervescente (oral); Yodo (EV) \\
\hline
TC / TCMS & Bario, Agua, Yodo, Aire y CO$_2$, Metilcelulosa, Manitol, Polietilenglicol \\
\hline
RM & Gadolinio (EV e intraarticular), Metilcelulosa, Manitol, Polietilenglicol \\
\hline
PET & Radiotrazadores: 18-FDG, Carbono-11, Flúor-18, Indio-111 (EV) \\
\hline
Ultrasonido (US) & Microburbujas (EV) \\
\hline
Angiografía (AD) & Yodo/Gadolinio y CO$_2$ (arterial); Yodo (EV) \\
\hline
\end{tabular}
\end{center}

% ============================================================
\section*{Medios de Contraste Iodados (MCI)}
% ============================================================

El \textbf{anillo de benceno} es el sustrato básico. Requiere \textbf{tres átomos de yodo en un monómero} (seis en un dímero) para lograr opacidad radiológica adecuada.

\begin{itemize}[noitemsep]
  \item C2, C4, C6: átomos de \textbf{yodo}.
  \item C1: grupo carboxilo (-COOH) o hidroxilo (-OH).
  \item C3 y C5: radicales orgánicos (determinan \textit{solubilidad}, \textit{unión a proteínas} y \textit{vía de excreción}).
\end{itemize}

Excreción predominantemente \textbf{renal}. Vida media $\approx$ 1 hora en individuos sanos.

\subsection*{Clasificación fisicoquímica}

\renewcommand{\arraystretch}{1.35}
\begin{center}
\begin{tabular}{p{2.8cm} >{\centering}p{2.5cm} >{\centering}p{2.5cm} >{\centering}p{2.5cm} >{\centering\arraybackslash}p{2.5cm}}
\hline
\textbf{Propiedad} & \textbf{\color{rojoriesgo}Alta OSM} & \textbf{Baja OSM} & \textbf{Baja OSM} & \textbf{\color{verdeclave}Isoosmolar} \\
\hline
OSM (mOsm/kg) & $\sim$2100 & 577 & 600--900 & 290 \\
\hline
Ionicidad & Iónico & Iónico & No iónico & No iónico \\
\hline
Estructura & Monómero & Dímero & Monómero & Dímero \\
\hline
Viscosidad 37°C & 8,4 & 9,5 & 7,8--11,2 & 11,1 \\
\hline
Ejemplo & Diatrizoato & Hexabrix & Omnipaque, Xenetix & Visipaque \\
\hline
\end{tabular}
\end{center}

\begin{cajakeypoints}
\textbf{Concepto clave:} Los contrastes \textbf{no iónicos, dímeros e isoosmolares} presentan \textbf{menor incidencia de reacciones adversas}. Esto no implica que sean los más comercializados.
\end{cajakeypoints}

\subsubsection*{Viscosidad}
Resistencia de los fluidos a la circulación. \textbf{Relación inversa} con temperatura y osmolaridad (se calientan a 37°C). \textbf{Relación directa} con tamaño molecular, lo que influye en la quimiotoxicidad.

% ============================================================
\section*{Gadolinio (Gd)}
% ============================================================

Metal \textbf{paramagnético} de la familia de los lantánidos. Acelera el tiempo de relajación (TR) en RM. \textbf{Aumenta la señal en T1} y la \textbf{disminuye en T2}. En forma libre es tóxico; se neutraliza con un \textbf{agente quelante}. Excreción \textbf{renal}.

\begin{itemize}[noitemsep]
  \item \textbf{Quelante lineal}: menor estabilidad, mayor riesgo de Gd libre (\textit{transmetilación}).
  \item \textbf{Quelante cíclico}: mayor estabilidad termodinámica. \textbf{Preferido clínicamente.}
\end{itemize}

% ============================================================
\section*{Reacciones Adversas}
% ============================================================

\subsection*{Por fisiopatología}

\subsubsection*{Fisicoquímicas o quimiotóxicas}
\textbf{Dosis-dependientes.} Acción directa sobre células, tejidos y sistemas enzimáticos. En vasos: daño endotelial, deformación de glóbulos rojos, trombosis.

\subsubsection*{Idiosincráticas o por hipersensibilidad}
\textbf{Independientes de la dosis. Impredecibles.}
\begin{itemize}[noitemsep]
  \item \textbf{Inmediata}: dentro de los \textbf{60 minutos} posteriores.
  \item \textbf{Tardía}: entre \textbf{3 horas y 2 días}; se autolimitan en 1 a 7 días.
\end{itemize}

\subsection*{Por gravedad}

\renewcommand{\arraystretch}{1.35}
\begin{center}
\begin{tabular}{>{\bfseries}p{2.2cm} p{1.8cm} p{4.5cm} p{5cm}}
\hline
Gravedad & Frec. & Manejo & Manifestaciones \\
\hline
\color{verdeclave}Leves & 98\% & Autolimitadas. Observación. & Náuseas, vómitos leves, mareos, exantema, congestión nasal. \\[0.5em]
\hline
\color{naranjadestacado}Moderadas & $\sim$1\% & Requieren tratamiento. & Broncoespasmo, disnea, taquicardia, urticaria extensa, hipotensión. \\[0.5em]
\hline
\color{rojoriesgo}Graves & $\sim$1\% & Tratamiento e internación. & Edema laríngeo, shock, paro cardiorrespiratorio, convulsiones. \\
\hline
\end{tabular}
\end{center}


\begin{cajariesgo}
\textbf{Importante:} El calor, el deseo de orinar y el gusto metálico son \textbf{respuestas fisiológicas normales}, no reacciones adversas leves.
\end{cajariesgo}

\subsection*{Reacciones adversas del Gadolinio}

\begin{cajariesgo}
\textbf{Fibrosis Sistémica Nefrogénica (FSN):} reacción adversa \textit{exclusiva} del gadolinio. Solo en pacientes con \textbf{insuficiencia renal grave}. Presenta dolor generalizado, prurito, eritema y edema de tejidos blandos.
\end{cajariesgo}

% ============================================================
\section*{Factores de Riesgo}
% ============================================================

\begin{itemize}[noitemsep]
  \item Reacción alérgica previa al MC (\textbf{5 veces} más riesgo).
  \item Asma bronquial (\textbf{6 a 10 veces} más riesgo).
  \item Cualquier tipo de alergia (\textbf{3 veces} más riesgo).
  \item Hipertiroidismo, ansiedad, nefrópatas, diabéticos, cardiópatas.
  \item \textbf{Contraindicaciones absolutas:} mieloma múltiple, feocromocitoma, deshidratación.
\end{itemize}

\textbf{Prevención:} interrogatorio, consentimiento informado, ayuno 4 horas, hidratación, creatinina en mayores de 60 años (umbral: $>$1,1 mg/dl en mujeres; $>$1,3 mg/dl en hombres).

\clearpage

% ============================================================
\section*{Nefropatía Inducida por Contraste (NIC)}
% ============================================================

\begin{cajadefinicion}
La \textbf{NIC} es un \textit{rápido deterioro de la función renal} luego de la administración de MC. Principal predictor: \textbf{disfunción renal previa}.
\end{cajadefinicion}

{\renewcommand{\arraystretch}{1.35}
\begin{center}
\begin{tabular}{>{\bfseries}p{3.5cm} p{3cm} p{6cm}}
\hline
FG (ml/min) & Riesgo & Conducta \\
\hline
$>$ 60 & Bajo & Sin profilaxis. \\
\hline
30 -- 60 & Moderado & Tratamiento preventivo. \\
\hline
$<$ 30 & Elevado & Tratamiento preventivo intensivo. \\
\hline
$<$ 15 & Máximo & Evaluación por nefrología; posible diálisis. \\
\hline
\end{tabular}
\end{center}


% ============================================================
\section*{Extravasación de MC}
% ============================================================

Salida del MC del vaso hacia los tejidos blandos. Síntomas: dolor, hinchazón, hematoma, ulceración, síndrome compartimental.

\textbf{Mecanismos de daño:} osmolaridad (a mayor OSM, mayor daño), citotoxicidad (mayor en iónicos), volumen (a mayor volumen, mayor daño), compresión mecánica.

\subsection*{Protocolo de actuación}
\begin{enumerate}[noitemsep]
  \item Retirar la vía.
  \item Estimar el volumen extravasado (remanente en la bomba).
  \item Detectar la zona y drenar. \textbf{No masajear.}
  \item Elevar el miembro afectado por encima del corazón.
  \item Aplicar hielo.
  \item Observar y documentar hasta mejoría.
  \item Si síntomas progresan luego de \textbf{2 a 4 horas}: consulta con cirujano.
\end{enumerate}

% ============================================================
\section*{Consentimiento Informado}
% ============================================================

El paciente \textbf{decide y autoriza de forma consciente, libre y responsable}. Requiere explicación oral y escrita, sin coacción. Debe ser impreso, sin abreviaturas ni enmendaduras, y firmado por el paciente o su representante legal.

Es un \textbf{deber del médico} y un \textbf{derecho del paciente}.

\clearpage

% ============================================================
\section*{Situaciones Especiales}
% ============================================================

\subsection*{Embarazo}
Gadolinio y MCI \textbf{atraviesan la barrera placentaria}. Evitar salvo necesidad absoluta. Riesgo conocido con MCI: \textbf{depresión de la función tiroidea fetal}.

\subsection*{Lactancia}
Menos del \textbf{1\%} de la dosis se excreta en leche materna; de ese porcentaje, menos del 1\% es absorbido por el lactante. Precaución: amamantar antes de la inyección y extraer leche previamente para el siguiente amamantamiento.

% ============================================================
\section*{Bomba Inyectora}
% ============================================================

Inyección segura y reproducible. Parámetros: volumen (ml), flujo (ml/s), presión (PSI), \textit{rise time} (s). Útil en TC, RM y AD.

\textbf{Secuencia de armado:}
\begin{enumerate}[noitemsep]
  \item Introducir vaso con jeringa en el cabezal.
  \item Subir el pistón hasta arriba.
  \item Colocar prolongador para rellenar.
  \item Rellenar contraste.
  \item Rellenar suero (dilución si monocabezal).
  \item Purgar el sistema.
  \item Orientar la bomba hacia abajo.
  \item Introducir los valores de inyección.
\end{enumerate}

% ============================================================
\section*{Esquema Mental --- Repaso Final}
% ============================================================

\begin{cajakeypoints}
\begin{center}
\textbf{\color{azulprincipal}
  MC $\rightarrow$ 3 clasificaciones $\rightarrow$ MCI: anillo bencénico $\rightarrow$ propiedades fisicoquímicas}\\[0.4em]
\textbf{\color{naranjadestacado}
  Mejor MC: no iónico + dímero + isoosmolar (Visipaque)}\\[0.3em]
\textbf{\color{rojoriesgo}
  RA: fisicoquímica (dosis) vs.\ idiosincrática (impredecible)}\\[0.3em]
Leves 98\% \quad Moderadas $\sim$1\% \quad Graves $\sim$1\%\\[0.3em]
\textbf{NIC}: FG $<$ 60 $\Rightarrow$ profilaxis \quad FG $<$ 15 $\Rightarrow$ nefrología\\[0.3em]
\textbf{Extravasación}: retirar vía $\to$ elevar $\to$ no masajear $\to$ hielo\\[0.3em]
\textbf{Gd}: paramagnético, quelante cíclico, FSN solo en insuf.\ renal grave
\end{center}
\end{cajakeypoints}

\vspace{1em}
\hrule
\vspace{0.5em}
{\footnotesize \textit{Fuente: material de clase --- Imagenología Especializada I, Lic. Prado y Lic. Canclini, UdelaR 2026.}}